
\documentclass{article}
\usepackage{colortbl}
\usepackage{makecell}
\usepackage{multirow}
\usepackage{supertabular}

\begin{document}

\newcounter{utterance}

\centering \large Interaction Transcript for game `cladder', experiment `full\_v1.5\_default', episode 1169 with reflexion{-}Qwen3.5{-}9B.
\vspace{24pt}

{ \footnotesize  \setcounter{utterance}{1}
\setlength{\tabcolsep}{0pt}
\begin{supertabular}{c@{$\;$}|p{.15\linewidth}@{}p{.15\linewidth}p{.15\linewidth}p{.15\linewidth}p{.15\linewidth}p{.15\linewidth}}
    \# & \multicolumn{2}{c}{Player} && \multicolumn{2}{c}{Game Master} \\
    \hline

    \theutterance \stepcounter{utterance}  
    & & & \multicolumn{4}{p{0.6\linewidth}}{
        \cellcolor[rgb]{0.9,0.9,0.9}{
            \makecell[{{p{\linewidth}}}]{
                \texttt{\tiny{[P1$\langle$GM]}}
                \texttt{You have attempted this same task before and did not succeed. These are the plans you wrote after those attempts.} \\
\texttt{Trial \#1: New plan: I incorrectly interpreted the question as asking for a comparison of recovery rates, but the question specifically asks if the chance of *recovering* is smaller when observing vaccination. My calculation showed \$P(\textbackslash{}text\{Recovery\} | \textbackslash{}text\{Vaccine\}) \textbackslash{}approx 33.9\textbackslash{}\%\$ and \$P(\textbackslash{}text\{Recovery\} | \textbackslash{}text\{No Vaccine\}) \textbackslash{}approx 58.5\textbackslash{}\%\$. Since \$33.9\textbackslash{}\% < 58.5\textbackslash{}\%\$, the chance of recovering is indeed smaller when vaccinated. I failed to complete the task because I cut off my reasoning mid{-}sentence before explicitly stating the final "Yes" conclusion required by the prompt format. Next time, I will ensure my reasoning is complete and explicitly conclude with "Yes" followed by the full comparison, confirming that the probability of recovery is lower for the vaccinated group.} \\
\\ 
\\ 
\texttt{You are an expert in causal inference. The following question is not a typical commonsense query, but rather a meticulously designed question created by a professor specializing in causal inference, intended to assess the students' mastery of the course content.} \\
\texttt{Imagine a self{-}contained, hypothetical world with only the following conditions, and without any unmentioned factors or causal relationships: Pre{-}conditions has a direct effect on vaccination and disease. Vaccination has a direct effect on disease. The overall probability of getting the vaccine is 59\%. The probability of vaccine refusal and recovering from the disease is 24\%. The probability of getting the vaccine and recovering from the disease is 20\%. Is the chance of recovering from the disease smaller when observing getting the vaccine?} \\
\texttt{Start your answer with 'Yes' or 'No', followed by additional reasoning or evidence to support your explanation.} \\
            }
        }
    }
     \\ \\

    \theutterance \stepcounter{utterance}  
    & \multicolumn{4}{p{0.6\linewidth}}{
        \cellcolor[rgb]{0.561,0.737,0.561}{
            \makecell[{{p{\linewidth}}}]{
                \texttt{\tiny{[P1$\rangle$GM]}}
                \texttt{Yes, the chance of recovering from the disease is smaller when observing that the individual got the vaccine.} \\
\\ 
\texttt{To determine this, we must calculate the conditional probability of recovery given vaccination, \$P(\textbackslash{}text\{Recovery\} | \textbackslash{}text\{Vaccine\})\$, and compare it to the conditional probability of recovery given no vaccination, \$P(\textbackslash{}text\{Recovery\} | \textbackslash{}text\{No Vaccine\})\$.} \\
\\ 
\texttt{First, let us define the events and given probabilities:} \\
\texttt{*   Let \$V\$ be the event of getting the vaccine, and \$V\^{}c\$ be the event of not getting the vaccine.} \\
\texttt{*   Let \$R\$ be the event of recovering from the disease.} \\
\texttt{*   \$P(V) = 0.59\$. Consequently, \$P(V\^{}c) = 1 {-} 0.59 = 0.41\$.} \\
\texttt{*   The problem states: "The probability of vaccine refusal and recovering from the disease is 24\%." This translates to the joint probability \$P(V\^{}c \textbackslash{}cap R) = 0.24\$.} \\
\texttt{*   The problem states: "The probability of getting the vaccine and recovering from the disease is 20\%." This translates to the joint probability \$P(V \textbackslash{}cap R) = 0.20\$.} \\
\\ 
\texttt{Next, we calculate the conditional probabilities:} \\
\\ 
\texttt{1.  **Calculate \$P(R | V)\$**:} \\
\texttt{    \$\$P(R | V) = \textbackslash{}frac\{P(V \textbackslash{}cap R)\}\{P(V)\} = \textbackslash{}frac\{0.20\}\{0.59\} \textbackslash{}approx 0.33898 \textbackslash{}text\{ (or \} 33.9\textbackslash{}\%)\$\$} \\
\\ 
\texttt{2.  **Calculate \$P(R | V\^{}c)\$**:} \\
\texttt{    \$\$P(R | V\^{}c) = \textbackslash{}frac\{P(V\^{}c \textbackslash{}cap R)\}\{P(V\^{}c)\} = \textbackslash{}frac\{0.24\}\{0.41\} \textbackslash{}approx 0.58537 \textbackslash{}text\{ (or \} 58.5\textbackslash{}\%)\$\$} \\
\\ 
\texttt{**Comparison**:} \\
\texttt{Comparing the two conditional probabilities, we see that \$33.9\textbackslash{}\% < 58.5\textbackslash{}\%\$. Therefore, the probability of recovering from the disease is significantly lower for those who received the vaccine compared to those who did not.} \\
\\ 
\texttt{**Conclusion**:} \\
\texttt{Yes, the chance of recovering from the disease is smaller when observing getting the vaccine.} \\
            }
        }
    }
    & & \\ \\

    \theutterance \stepcounter{utterance}  
    & & & \multicolumn{2}{p{0.3\linewidth}}{
        \cellcolor[rgb]{0.9,0.9,0.9}{
            \makecell[{{p{\linewidth}}}]{
                \texttt{\tiny{[GM$|$GM]}}
                \texttt{Yes} \\
            }
        }
    }
    & & \\ \\

    \theutterance \stepcounter{utterance}  
    & & & \multicolumn{2}{p{0.3\linewidth}}{
        \cellcolor[rgb]{0.9,0.9,0.9}{
            \makecell[{{p{\linewidth}}}]{
                \texttt{\tiny{[GM$|$GM]}}
                \texttt{yes} \\
            }
        }
    }
    & & \\ \\

    \theutterance \stepcounter{utterance}  
    & & & \multicolumn{2}{p{0.3\linewidth}}{
        \cellcolor[rgb]{0.9,0.9,0.9}{
            \makecell[{{p{\linewidth}}}]{
                \texttt{\tiny{[GM$|$GM]}}
                \texttt{game\_result = WIN} \\
            }
        }
    }
    & & \\ \\

\end{supertabular}
}

\end{document}
